% ============================================================
%  MTH636 – Game Theory | Lecture 25
%  Kuhn's Theorem & Equilibria in Behavior Strategies
% ============================================================
\documentclass[12pt, a4paper]{article}

% ============================================================
%  MTH636 -- Game Theory
%  Notes Preamble — Color-Coded, TikZ-Enhanced
% ============================================================

% --- Core packages ---
\usepackage[T1]{fontenc}
\usepackage[utf8]{inputenc}
\usepackage{lmodern}
\usepackage[margin=2.5cm, top=2.8cm, bottom=2.8cm]{geometry}
\usepackage{microtype}

% --- Math ---
\usepackage{amsmath, amssymb, amsthm, mathtools}
\usepackage{bm}        % bold math
\usepackage{cancel}    % strike-through in math

% --- Color & Boxes ---
\usepackage[dvipsnames, svgnames, x11names, table]{xcolor}
\usepackage[most]{tcolorbox}
\tcbuselibrary{breakable, skins, theorems}

% --- TikZ & Diagrams ---
\usepackage{tikz}
\usepackage{tikz-qtree}
\usetikzlibrary{arrows.meta, positioning, shapes.geometric, fit, backgrounds,
                decorations.pathreplacing, calc, matrix, trees}
\usepackage{pgfplots}
\pgfplotsset{compat=1.18}

% --- Tables ---
\usepackage{booktabs}
\usepackage{multirow}
\usepackage{array}
\usepackage{tabularx}
\usepackage{makecell}

% --- Typography & Layout ---
\usepackage{enumitem}
\usepackage{parskip}
\usepackage{fancyhdr}
\usepackage{titlesec}
\usepackage{titletoc}
\usepackage{hyperref}
\usepackage{graphicx}
\usepackage{float}
\usepackage{caption}

% --- Color Palette ---
\definecolor{defblue}{HTML}{1A5276}      % Definition header
\definecolor{defbluebg}{HTML}{D6EAF8}    % Definition background
\definecolor{thmgreen}{HTML}{1E8449}     % Theorem header
\definecolor{thmbg}{HTML}{D5F5E3}        % Theorem background
\definecolor{exorange}{HTML}{B7770D}     % Example header
\definecolor{exbg}{HTML}{FEF9E7}         % Example background
\definecolor{warnred}{HTML}{922B21}      % Warning/Key Insight header
\definecolor{warnbg}{HTML}{FADBD8}       % Warning background
\definecolor{noticegray}{HTML}{2C3E50}   % Note header
\definecolor{noticebg}{HTML}{EBF5FB}     % Note background
\definecolor{taketawayyellow}{HTML}{7D6608} % Key Takeaway header
\definecolor{takeawaybg}{HTML}{FDFEFE}   % Takeaway background
\definecolor{accentpurple}{HTML}{6C3483} % Accent / proof
\definecolor{lightgray}{HTML}{F2F3F4}    % shading

% ============================================================
%  TCOLORBOX ENVIRONMENTS
% ============================================================

% -- Definition (Blue) --
\newtcolorbox{defbox}[2][]{
  enhanced, breakable,
  colback=defbluebg, colframe=defblue, coltitle=white,
  fonttitle=\bfseries\sffamily,
  title={Definition: #2},
  attach boxed title to top left={yshift=-2mm, xshift=4mm},
  boxed title style={colback=defblue, rounded corners},
  arc=4pt, boxrule=1.2pt,
  left=6pt, right=6pt, top=8pt, bottom=6pt,
  #1
}

% -- Theorem / Proposition (Green) --
\newtcolorbox{thmbox}[2][]{
  enhanced, breakable,
  colback=thmbg, colframe=thmgreen, coltitle=white,
  fonttitle=\bfseries\sffamily,
  title={Theorem: #2},
  attach boxed title to top left={yshift=-2mm, xshift=4mm},
  boxed title style={colback=thmgreen, rounded corners},
  arc=4pt, boxrule=1.2pt,
  left=6pt, right=6pt, top=8pt, bottom=6pt,
  #1
}

% -- Proposition (Green variant) --
\newtcolorbox{propbox}[2][]{
  enhanced, breakable,
  colback=thmbg, colframe=thmgreen!70!black, coltitle=white,
  fonttitle=\bfseries\sffamily,
  title={Proposition: #2},
  attach boxed title to top left={yshift=-2mm, xshift=4mm},
  boxed title style={colback=thmgreen!70!black, rounded corners},
  arc=4pt, boxrule=1.2pt,
  left=6pt, right=6pt, top=8pt, bottom=6pt,
  #1
}

% -- Example (Orange/Yellow) --
\newtcolorbox{exbox}[2][]{
  enhanced, breakable,
  colback=exbg, colframe=exorange, coltitle=white,
  fonttitle=\bfseries\sffamily,
  title={Example: #2},
  attach boxed title to top left={yshift=-2mm, xshift=4mm},
  boxed title style={colback=exorange, rounded corners},
  arc=4pt, boxrule=1.2pt,
  left=6pt, right=6pt, top=8pt, bottom=6pt,
  #1
}

% -- Key Insight / Warning (Red) --
\newtcolorbox{warnbox}[2][]{
  enhanced, breakable,
  colback=warnbg, colframe=warnred, coltitle=white,
  fonttitle=\bfseries\sffamily,
  title={Key Insight: #2},
  attach boxed title to top left={yshift=-2mm, xshift=4mm},
  boxed title style={colback=warnred, rounded corners},
  arc=4pt, boxrule=1.2pt,
  left=6pt, right=6pt, top=8pt, bottom=6pt,
  #1
}

% -- Note (Gray-Blue) --
\newtcolorbox{notebox}[1][]{
  enhanced, breakable,
  colback=noticebg, colframe=noticegray!60, coltitle=white,
  fonttitle=\bfseries\sffamily,
  title={Note},
  attach boxed title to top left={yshift=-2mm, xshift=4mm},
  boxed title style={colback=noticegray!60, rounded corners},
  arc=4pt, boxrule=0.8pt,
  left=6pt, right=6pt, top=8pt, bottom=6pt,
  #1
}

% -- Key Takeaway --
\newtcolorbox{takebox}[1][]{
  enhanced, breakable,
  colback=takeawaybg, colframe=taketawayyellow, coltitle=white,
  fonttitle=\bfseries\sffamily,
  title={Key Takeaway},
  attach boxed title to top left={yshift=-2mm, xshift=4mm},
  boxed title style={colback=taketawayyellow, rounded corners},
  arc=4pt, boxrule=1.2pt,
  left=6pt, right=6pt, top=8pt, bottom=6pt,
  #1
}

% -- Intuition --
\newtcolorbox{intubox}[1][]{
  enhanced, breakable,
  colback=accentpurple!8, colframe=accentpurple!60, coltitle=white,
  fonttitle=\bfseries\sffamily,
  title={Intuition},
  attach boxed title to top left={yshift=-2mm, xshift=4mm},
  boxed title style={colback=accentpurple!80, rounded corners},
  arc=4pt, boxrule=0.8pt,
  left=6pt, right=6pt, top=8pt, bottom=6pt,
  #1
}

% ============================================================
%  PAYOFF MATRIX HELPER
%  Usage: inside tabular, use \payoff{a}{b} for (a,b)
% ============================================================
\newcommand{\payoff}[2]{$(#1,\;#2)$}
\newcommand{\nash}[2]{\underline{$#1$},\underline{$#2$}}  % underline NE payoffs
\newcommand{\Nset}{\mathbb{N}}
\newcommand{\Rset}{\mathbb{R}}
\newcommand{\Sset}{\mathcal{S}}
\newcommand{\Ui}{u_i}

% ============================================================
%  SECTION FORMATTING
% ============================================================
\titleformat{\section}
  {\large\bfseries\sffamily\color{defblue}}
  {\thesection}{1em}{}[\titlerule]

\titleformat{\subsection}
  {\normalsize\bfseries\sffamily\color{thmgreen!80!black}}
  {\thesubsection}{1em}{}

\titleformat{\subsubsection}
  {\normalsize\itshape\sffamily\color{exorange!80!black}}
  {\thesubsubsection}{1em}{}

% ============================================================
%  HEADER / FOOTER
% ============================================================
\pagestyle{fancy}
\fancyhf{}
\fancyhead[L]{\sffamily\small\color{defblue} MTH636 -- Game Theory}
\fancyhead[R]{\sffamily\small\color{defblue} \leftmark}
\fancyfoot[C]{\sffamily\small\thepage}
\renewcommand{\headrulewidth}{0.4pt}

% ============================================================
%  HYPERREF SETUP
% ============================================================
\hypersetup{
  colorlinks=true,
  linkcolor=defblue,
  urlcolor=thmgreen,
  citecolor=accentpurple,
  bookmarksnumbered=true,
  pdftitle={MTH636 Game Theory Notes},
  pdfauthor={MTH636 Student Notes}
}

% ============================================================
%  MISC HELPERS
% ============================================================
\newcommand{\BR}{\text{BR}}        % Best Response
\newcommand{\NE}{\text{NE}}        % Nash Equilibrium
\newcommand{\SPNE}{\text{SPNE}}    % Subgame Perfect NE
\newcommand{\Dom}{\text{dom}}      % Dominance
\newcommand{\mixs}[1]{\sigma_{#1}} % Mixed strategy
\newcommand{\pures}[1]{s_{#1}}     % Pure strategy


% ============================================================
\begin{document}

% Title
\begin{center}
  {\LARGE\bfseries\sffamily\color{defblue} MTH636 -- Game Theory}\\[6pt]
  {\Large\sffamily Lecture 25: Kuhn's Theorem \& Equilibria in Behavior Strategies}\\[4pt]
\end{center}
\vspace{0.5cm}
\hrule height 1.5pt
\vspace{0.3cm}

% ============================================================
\section{Recap: Constructing an Equivalent Behavior Strategy from a Mixed Strategy}

From the previous lecture, when we have a mixed strategy $\sigma_1$ for Player I and want to find an equivalent behavior strategy $b_1$, we use the following approach via a worked example.

\subsection{Example: Two-Level Game Tree}

Consider a game tree where Player I has two information sets: $h_1 = \{n_0\}$ (the root, actions $L$ with prob.\ $u$ and $R$ with prob.\ $1-u$) and $h_2 = \{n_1, n_2\}$ (actions $\ell$ and $r$). Player I's pure strategies are $L\ell$, $Lr$, $R\ell$, $Rr$.

Suppose the mixed strategy satisfies:
\[
  \sigma_1(L\ell) = \sigma_1(Rr) = \tfrac{1}{2}.
\]
The terminal nodes have labels as shown; for Player II the behavior strategy at $h_2$ uses probabilities $y$ and $(1-y)$ for $\ell$ and $r$:

\begin{center}
  \begin{tabular}{ll}
    $\sigma_1(L\ell) = \tfrac{1}{2}$, & $\sigma_1(Rr) = \tfrac{1}{2}$.
  \end{tabular}
\end{center}

Let $b_1(L;\, n_0) = u$ and $b_1(\ell;\, h_2) = y$. Then the reaching probabilities give:
\[
  uy = \tfrac{1}{2}, \quad u(1-y) = 0, \quad (1-u)y = 0, \quad (1-u)(1-y) = \tfrac{1}{2}.
\]
This system has the solution: $u = \tfrac{1}{2}$, $y = 1$ is \emph{not} consistent; instead the only solution is $u=1, y=1/2$ fails too. In fact the compatible solution is $u = 1/2$, which implies from $u(1-y)=0$ that $y=1$, contradiction with $uy=1/2$ giving $y=1$. Re-examining: $uy = 1/2$ and $(1-u)(1-y) = 1/2$, and $u(1-y)=0$ so $u=0$ or $y=1$. Since $u=0$ contradicts $uy=1/2$, we need $y=1$. But then $(1-u)(1-y)=0\neq 1/2$. Hence \textbf{no solution exists} for this imperfect-recall tree.

\medskip
Now consider a different (perfect-recall) tree. Suppose:
\[
  \sigma_1(AL) = \tfrac{1}{5}, \quad \sigma_1(Ar) = \tfrac{1}{10}, \quad
  \sigma_1(BL) = \tfrac{2}{5}, \quad \sigma_1(Br) = \tfrac{3}{10}.
\]
Here $\sigma_1(A\cdot) = \tfrac{1}{5}+\tfrac{1}{10} = \tfrac{3}{10}$ and $\sigma_1(B\cdot)= \tfrac{2}{5}+\tfrac{3}{10}=\tfrac{7}{10}$... (see boxed equations below).

For the tree with $n_0$ (Player I, actions $A$ prob.\ $u$ and $B$ prob.\ $1-u$), $n_1$ (Player II, actions $L$ and $R$), and an information set $h_2 = \{n_3, n_4\}$ (Player I, actions $\ell$ prob.\ $y$ and $r$ prob.\ $1-y$):
\[
  \boxed{ny = \tfrac{1}{5},\quad n(1-y) = \tfrac{1}{10},\quad (1-n)y = \tfrac{2}{5},\quad (1-n)(1-y) = \tfrac{3}{10}.}
\]
Solving: $n(y + 1-y) = ny + n(1-y) = \tfrac{1}{5}+\tfrac{1}{10} = \tfrac{3}{10}$, so $n = \tfrac{3}{10}$... Checking sign: from $ny = \tfrac{1}{5}$ with $n = 3/10$: $y = \tfrac{1/5}{3/10} = \tfrac{2}{3}$. Similarly $(1-n) = 7/10$.

% ============================================================
\section{Constructing $b_i$ from $\sigma_i$ -- General Formula}

\begin{defbox}{Converting a Mixed Strategy to a Behavior Strategy}
  Let $S_i(h_2) \subseteq S_i$ denote the set of strategies of Player $i$ such that the game \textbf{possibly reaches} $h_2$ (i.e., the information set $h_2$ is reachable under $s_i$).

  \medskip
  More precisely, for information set $h_2$, define:
  \[
    S_i(\ell;\,h_2) \subseteq S_i(h_2)
  \]
  as the set of strategies in $S_i(h_2)$ that \textbf{choose action $\ell$ at $h_2$}.

  \medskip
  Also define:
  \begin{align*}
    \sigma_i\!\bigl(S_i(h_2)\bigr)   &: \text{ the total probability of } S_i(h_2) \text{ under } \sigma_i, \\
    \sigma_i\!\bigl(S_i(\ell;\,h_2)\bigr) &: \text{ the total probability of } S_i(\ell;\,h_2) \text{ under } \sigma_i.
  \end{align*}

  The behavior strategy $b_i$ is then constructed as:
  \[
    b_i(\ell;\,h_2) =
    \begin{cases}
      \dfrac{\sigma_i\!\bigl(S_i(\ell;\,h_2)\bigr)}{\sigma_i\!\bigl(S_i(h_2)\bigr)} & \text{provided } \sigma_i\!\bigl(S_i(h_2)\bigr) > 0, \\[8pt]
      \alpha \in [0,1] & \text{o.f.\ (otherwise free)},
    \end{cases}
  \]
  and similarly:
  \[
    b_i(r;\,h_2) =
    \begin{cases}
      \dfrac{\sigma_i\!\bigl(S_i(r;\,h_2)\bigr)}{\sigma_i\!\bigl(S_i(h_2)\bigr)} & \text{provided } \sigma_i\!\bigl(S_i(h_2)\bigr) > 0, \\[8pt]
      1 - \alpha & \text{o.f.}
    \end{cases}
  \]
\end{defbox}

\subsection{Computed Example}

For the information set $h_1 = S_i$ (the root):
\[
  S_i(h_1) = S_i.
\]
\[
  S_i(A;\,h_1) = \{A\ell,\, Ar\}.
\]
\[
  b_i(A;\,h_1)
  = \frac{\sigma_i\!\bigl(S_i(A;\,h_1)\bigr)}{\sigma_i(S_i)}
  = \frac{\sigma_i(A\ell) + \sigma_i(Ar)}{1}
  = \tfrac{1}{5} + \tfrac{1}{10} = \tfrac{3}{10}.
\]

\begin{notebox}
  When $\sigma_i(S_i(h_2)) = 0$ (the information set is reached with zero probability under $\sigma_i$), the value of $b_i(\ell;\,h_2)$ can be chosen \emph{freely} (any $\alpha \in [0,1]$). This is why the conversion from mixed strategy to behavior strategy is generally non-unique only when some information sets are off the equilibrium path.
\end{notebox}

% ============================================================
\section{Kuhn's Theorem (1957)}

\begin{thmbox}{Kuhn's Theorem (1957)}
  In every game in extensive form, if Player $i$ has \textbf{perfect recall}, then for every mixed strategy of Player $i$, there exists a behavior strategy that is \textbf{equivalent} to the mixed strategy.
\end{thmbox}

\begin{notebox}
  Together with the result from Lecture 24 (that given any behavior strategy one can always construct an equivalent mixed strategy), Kuhn's Theorem establishes that \textbf{under perfect recall}, the sets of mixed strategies and behavior strategies are \emph{strategically equivalent} -- every mixed strategy has an equivalent behavior strategy and vice versa.
\end{notebox}

% ============================================================
\section{Equilibria in Behavior Strategies}

\begin{thmbox}{Existence of NE in Behavior Strategies}
  If all players in an extensive form game have \textbf{perfect recall}, then the game has a Nash Equilibrium in behavior strategies.
\end{thmbox}

\begin{proof}
  From the extensive form game, construct the corresponding strategic (normal) form game.

  Since the extensive form game is finite, the corresponding strategic form game is finite too.

  Now, the strategic form game has a mixed strategy Nash Equilibrium -- say $\sigma^* = (\sigma_1^*,\ldots,\sigma_n^*)$. For each $\sigma_i^*$, construct $b_i^*$, an outcome-equivalent behavior strategy (which exists by Kuhn's Theorem, since each player has perfect recall). We claim $b^* = (b_1^*,\ldots,b_n^*)$ is an NE in behavior strategies.

  By Remark R-2 from Lecture 24:
  \[
    U_i(\sigma^*) = U_i(b^*) \quad \forall\, i \in N.
  \]

  Assume for contradiction that $b^*$ is \textbf{not} an NE. Then $\exists\, k \in N$ such that:
  \[
    U_k(b^*) < U_k(b_k,\, b^*_{-k})
    \quad \text{for some behavior strategy } b_k.
  \]
  Now, since $b_k$ is a behavior strategy, there exists an equivalent mixed strategy $\sigma_k$ (this direction always holds). Therefore:
  \[
    U_k(b^*) < U_k(b_k,\, b^*_{-k}) = U_k(\sigma_k,\, b^*_{-k}).
  \]
  Since $\sigma^*$ is equivalent to $b^*$ component-wise and $\sigma_k$ is equivalent to $b_k$:
  \[
    = U_k(\sigma_k,\, \sigma^*_{-k}) \leq U_k(\sigma^*), \quad \text{a contradiction.}
  \]
  (The last inequality holds because $\sigma^*$ is a Nash Equilibrium of the strategic form game, so no player can profitably deviate.)

  Hence $b^*$ is an NE in behavior strategies.
\end{proof}

\begin{takebox}
  \textbf{Key Takeaway.} For extensive form games with perfect recall:
  \begin{enumerate}
    \item Every mixed strategy $\Leftrightarrow$ some behavior strategy (Kuhn's Theorem, plus the easy direction).
    \item A Nash Equilibrium in \emph{mixed} strategies always induces a Nash Equilibrium in \emph{behavior} strategies.
    \item Behavior strategies are preferred in practice because they are far more compact to describe.
  \end{enumerate}
\end{takebox}

\end{document}
